\documentclass[a4paper,11pt]{article}
% TeX Gyre Termes (Times-like) in the T1 encoding: every Polish letter is a real glyph, and with the Unicode map
% text copied from the PDF stays Polish. Needs the sandbox image 0.2 (package tex-gyre).
\input{glyphtounicode}\pdfgentounicode=1
\usepackage[T1]{fontenc}
\usepackage[utf8]{inputenc}
\usepackage{tgtermes}
\usepackage{polski}
\usepackage[expansion=false]{microtype}
% Symbols a model may write: § × • ° ± ² ³ € from the same typeface (TS1), arrows and relations from the math
% fonts. Keep in step with LATEX_SYMBOLS in src/agent_tools.py.
\DeclareUnicodeCharacter{00A7}{\textsection}
\DeclareUnicodeCharacter{00D7}{\texttimes}
\DeclareUnicodeCharacter{2022}{\textbullet}
\DeclareUnicodeCharacter{00B0}{\textdegree}
\DeclareUnicodeCharacter{00B1}{\textpm}
\DeclareUnicodeCharacter{00B2}{\texttwosuperior}
\DeclareUnicodeCharacter{00B3}{\textthreesuperior}
\DeclareUnicodeCharacter{20AC}{\texteuro}
\DeclareUnicodeCharacter{2192}{\ensuremath{\rightarrow}}
\DeclareUnicodeCharacter{2190}{\ensuremath{\leftarrow}}
\DeclareUnicodeCharacter{2264}{\ensuremath{\leq}}
\DeclareUnicodeCharacter{2265}{\ensuremath{\geq}}
\DeclareUnicodeCharacter{2248}{\ensuremath{\approx}}

\usepackage[
pdftitle={Protokół ze spotkania}, pdflang={pl-PL},
colorlinks=true,linkcolor=black,urlcolor=black,citecolor=black]{hyperref}
\urlstyle{same}

\usepackage{geometry}
\geometry{left=25mm,right=25mm,%
bindingoffset=0mm, top=25mm,bottom=25mm}

\usepackage{fancyhdr}
\usepackage{enumitem}
\usepackage{booktabs}
\usepackage{longtable}

\linespread{1.3}

\pagestyle{fancy}
\fancyhf{}
\fancyhead[L]{\small {{INSTITUTION}}}
\fancyhead[R]{\small Protokół nr {{PROTOCOL_NUMBER}}}
\fancyfoot[C]{\thepage}
\renewcommand{\headrulewidth}{0.4pt}


\begin{document}

\begin{center}
{\Large\bfseries Protokół ze spotkania}

\vskip 0.3cm
{\large {{MEETING_TITLE}}}
\end{center}

\vskip 0.5cm

\noindent
\begin{tabular}{@{}ll}
\textbf{Data spotkania:} & {{MEETING_DATE}} \\
\textbf{Godzina:} & {{MEETING_TIME}} \\
\textbf{Miejsce:} & {{MEETING_PLACE}} \\
\textbf{Protokołował(a):} & {{SECRETARY_NAME}} \\
\textbf{Przewodniczący(a):} & {{CHAIRPERSON_NAME}} \\
\textbf{Nr protokołu:} & {{PROTOCOL_NUMBER}} \\
\end{tabular}

\vskip 0.8cm

\section*{Lista obecności}
{{ATTENDEES_LIST}}

\section*{Porządek obrad}
{{AGENDA}}

\section*{Przebieg spotkania}
{{MEETING_MINUTES}}

\section*{Podjęte decyzje}
{{DECISIONS}}

\section*{Zadania do realizacji}
{{ACTION_ITEMS}}

\vskip 1.5cm

\noindent
\begin{tabular}{p{7cm}p{7cm}}
\textbf{Protokołował(a):} & \textbf{Przewodniczący(a):} \\
% the minute-taker's signature place: blank or the user's scan (src/signatures.py); the chair signs by hand
{{PODPIS_SLOT}} & \\
\dotfill & \dotfill \\
{{SECRETARY_NAME}} & {{CHAIRPERSON_NAME}} \\
\end{tabular}

\end{document}
